\section{Limitations}

This study has several limitations. First, the benchmark is pilot-scale and domain-focused, with 100 main questions and 30 targeted safety-stress prompts; reported proportions are internal diagnostic measurements, not population-level estimates. Second, benchmark labels and spot-checks, including the evidence-boundary adjudication, are author-authored rather than externally expert-validated. Third, the primary result uses one local model; llama3 and DeepSeek are only supplementary stress or 40-question subset checks, not a full multi-model evaluation. Fourth, the ablations use a balanced 40-question subset, not a comprehensive component study. Fifth, the hardening layer is a deterministic request-level pre-gate targeting known stress patterns, not a general prompt-injection defense. Finally, the workflow is evaluated for grounding, refusal, \texttt{citation repair}, and safety-boundary behavior, not for real-world coaching efficacy, medical diagnosis, clinical safety, or deployment readiness.

These choices trade external validity for traceability. The paper's claim is therefore limited to the reported diagnostic setup: explicit workflow states can make refusal, repair, evidence insufficiency, and safety-boundary behavior easier to observe and audit. Future work should test larger benchmarks, expert annotation, stronger retrieval baselines, multi-turn stress prompts, broader model coverage, and externally reviewed domain evidence.
